\section{各 Case における結果の整理}
\label{sec:discussion_summary}

まず，提案手法を適用したシミュレーション結果を Case ごとに整理する．Case1 および Case2 では，床の沈み込み量が小さく，DS から SS への遷移時においても自重に起因する転倒モーメントが小さい条件であった．このため，ZMP，CoM 高さ，支持脚トルクのいずれにおいても，従来手法と提案手法の間に顕著な差は見られなかった．これらの条件では，もともと歩行の安定性が高く，提案手法による補正効果が顕在化しにくかったと考えられる．

一方，Case3 では，床の沈み込みが急激かつ大きく生じる条件であり，提案手法を適用した場合であっても転倒を回避することができなかった．この結果から，Case3 においては，提案手法による CoM 高さの制御が作用する以前に，沈み込みに起因する姿勢の不安定化が支配的となっており，重心高さの調整によって転倒挙動を緩和できる範囲を超えていることが示唆される．すなわち，本研究で提案する手法には適用可能な条件範囲が存在し，沈み込み量や沈み込み速度が過度に大きい場合には十分な効果を発揮できないことが確認された．

これに対して Case4 では，沈み込みが顕著に生じる条件であるものの，提案手法を適用することで歩行を継続することが可能となった．Case4 においては，DS から SS へ移行する直前および移行直後において CoM 高さが意図的に高く保たれ，その結果として横方向の倒れ込みが抑制された．これにより，SS 期間における ZMP の発散が防がれ，支持脚足裏領域内に ZMP を維持することが可能となったと考えられる．

さらに，Case4 では支持脚ロール方向トルクのピーク値も低減しており，CoM 高さの操作が力学的負荷の軽減にも寄与していることが確認された．以上の結果から，提案手法は，沈み込みによる不安定化が一定範囲内に収まる条件において，DS--SS 遷移時の転倒挙動を効果的に緩和し，歩行安定性を向上させる手法であるといえる．
